\documentclass[12pt,a4paper]{article}
\usepackage[utf8]{inputenc}
\usepackage[T2A]{fontenc}
\usepackage[russian,english]{babel}
\usepackage{geometry,booktabs,float,longtable,array,titlesec,fancyhdr,hyperref}
\newcolumntype{L}[1]{>{\raggedright\arraybackslash}p{#1}}
\geometry{left=2.5cm,right=1.5cm,top=2cm,bottom=2cm}
\hypersetup{colorlinks=true,urlcolor=blue,linkcolor=black}
\setlength{\emergencystretch}{2em}
\setlength{\headheight}{15pt}
\titleformat{\section}{\Large\bfseries}{\thesection}{1em}{}
\titleformat{\subsection}{\large\bfseries}{\thesubsection}{1em}{}
\pagestyle{fancy}
\fancyhf{}
\fancyhead[L]{Лабораторная работа №3: RAG-система}
\fancyhead[R]{\thepage}
\begin{document}
\selectlanguage{russian}
\hypersetup{pageanchor=false}
\begin{titlepage}
\begin{center}
\textbf{МИНИСТЕРСТВО НАУКИ И ВЫСШЕГО ОБРАЗОВАНИЯ РОССИЙСКОЙ ФЕДЕРАЦИИ}
\vspace{0.5cm}

\textbf{Федерального государственного автономного образовательного учреждения \\ высшего образования}
\vspace{0.3cm}

\textbf{«НАЦИОНАЛЬНЫЙ ИССЛЕДОВАТЕЛЬСКИЙ ТЕХНОЛОГИЧЕСКИЙ УНИВЕРСИТЕТ МИСИС»}
\vspace{2cm}

\textbf{Институт компьютерных наук}
\vspace{0.3cm}

\textbf{Кафедра Инженерной кибернетики}
\vspace{1cm}

\textbf{Лабораторная работа №3}
\vspace{0.3cm}

\textbf{по дисциплине «Современные интеллектуальные сетевые сервисы»}
\vspace{1.5cm}

\LARGE{\textbf{Разработка и исследование \\ RAG-системы для работы \\ с PDF документом}}
\vspace{2cm}
\end{center}
\begin{flushright}
\textbf{Выполнил:} \\
Студент группы МПИ-25-1-1 \\
Чун Донг Вук \\
\vspace{0.5cm}
\textbf{Проверил:} \\
Преподаватель \\
Корнет Мария Евгеньевна \\
\vspace{2cm}
\textbf{Оценка:} \hrulefill
\end{flushright}
\begin{center}\vfill Москва, 2026\end{center}
\end{titlepage}
\hypersetup{pageanchor=true}
\tableofcontents\newpage
\section{Цель и исходные данные}
Цель работы - разработать RAG-систему для ответов по PDF и исследовать влияние параметров поиска, промпта и памяти диалога. Исследуется статья OpenAI Training language models to follow instructions with human feedback, 68 страниц. Репозиторий проекта: \url{https://github.com/skyaker/aiservices}. Ноутбук: \url{https://github.com/skyaker/aiservices/blob/main/lab3_rag/rag_experiments.ipynb}.

Задачи работы: подготовить документ и индекс; сравнить полноту индекса, размеры и перекрытия фрагментов; сравнить лексический и векторный поиск; проверить ссылки, отказы модели и работу с уточнениями.

Использованы \texttt{llama3.2:1b} для генерации и \texttt{embeddinggemma} для эмбеддингов через Ollama. В отдельной проверке ссылок дополнительно использована \texttt{llama3.2:3b}; условия внутри каждого сравнения оставались одинаковыми. Страницы загружаются \texttt{PyPDFLoader}, делятся \texttt{Recursive\-Character\-Text\-Splitter} и индексируются \texttt{InMemory\-VectorStore}. Базовые параметры: 1000 символов, перекрытие 200, $k=3$, температура 0. Каждый фрагмент содержит идентификатор, источник и номер страницы.

В окружении Nix: LangChain 1.4.2, langchain-community 0.4.2, langchain-ollama 1.1.0, langchain-text-splitters 1.1.3, pypdf 6.21.4, requests 2.34.2. Файл \texttt{flake.lock} фиксирует версию набора пакетов.

Автоматическая оценка считает успешным поиск, в котором top-3 содержит фрагмент с заранее отмеченной страницей ответа: $H/N$. Это приближение к Hit Rate@3, а не полная проверка содержимого: фрагмент с нужным ответом может быть на другой странице, а совпадение страницы не гарантирует достаточность текста. Например, PPO-ptx описан на странице 9, хотя эталоном указана страница 4. Ответы модели оцениваются вручную: 0 - неверный или неподтвержденный, 1 - частично верный, 2 - верный и подтвержденный контекстом.

\section{Результаты}
\subsection{Подготовка и методика}
RAG сочетает поиск по документу с генерацией ответа по найденному контексту. В работе текст каждой страницы PDF разбивается отдельно; индекс хранится в памяти и пересоздается при изменении состава документов или параметров разбиения. Каждый запрос выполняет новый поиск трех фрагментов. Бинарный мешок слов сопоставляет множества слов, а плотный поиск сравнивает эмбеддинги модели. Позиция \texttt{start\_index} относится к извлеченному тексту конкретной страницы; номера в отчетных таблицах - физические номера страниц PDF, а не обязательно напечатанные в статье номера.

\subsection{Контрольные вопросы и эталоны}
Номера относятся к физическим страницам PDF. Одно и то же утверждение может повторяться в разных разделах, поэтому при ручной проверке учитывался текст найденного фрагмента, а не только указанная страница.
{\small\setlength{\tabcolsep}{4pt}\begin{longtable}{L{6.8cm}L{6.8cm}}
\caption{Вопросы и эталонные ответы}\\\toprule
Вопрос & Эталонный ответ и страница\\\midrule\endfirsthead
\toprule Вопрос & Эталонный ответ и страница\\\midrule\endhead
Три стадии обучения InstructGPT? & SFT, обучение reward model, PPO; стр. 3: supervised fine-tuning, reward model training, reinforcement learning.\\
Зачем разметчики писали первые промпты? & Для первоначального набора инструкций; стр. 7: needed an initial source of instruction-like prompts.\\
Что добавляет PPO-ptx? & Обновления по данным предварительного обучения; стр. 9: mixing the pretraining gradients into the PPO gradients.\\
Какая меньшая модель предпочтена GPT-3 175B? & InstructGPT 1.3B; стр. 1: 1.3B parameter InstructGPT model are preferred.\\
Какие три свойства поведения требовались? & Полезность, правдивость, безопасность; стр. 2: helpful, honest, and harmless.\\
Наиболее частая категория запросов API? & Generation, 45.6\%; стр. 6: Generation 45.6\%.\\
Где набирали разметчиков? & Upwork и Scale AI; стр. 18: hired via Upwork or Scale AI.\\
Какой тест используют для токсичности? & RealToxicityPrompts; стр. 10: human evaluations of toxicity on the RealToxicityPrompts dataset.\\
Почему набор ценностей ограничен? & Оценки зависят от группы разметчиков; стр. 19: not representative of the full spectrum of people.\\
Что делали с равными сравнениями? & Исключали из данных; стр. 42: Ties were dropped.\\\bottomrule
\end{longtable}}
\subsection{Полнота индекса}
Проверены три вопроса по началу статьи и три по последующим разделам. В первых пяти фрагментах нет полного описания стадий обучения и сведений о наборе разметчиков, поэтому промпт не может вернуть отсутствующие данные в контекст. При этом определение отличия PPO-ptx от PPO присутствует в подписи к рисунку на странице 2.
\begin{table}[H]\centering\caption{Полнота индекса: проверка содержания}\begin{tabular}{lcccc}\toprule Индекс&Начало, из 3&Далее, из 3&Hit Rate@3&Баллы 0--2\\\midrule Первые 5&1&0&0.17&0.00\\ Полный&3&2&0.83&0.00\\\bottomrule\end{tabular}\end{table}
Ручная проверка выявила ответы о PPO-ptx во фрагменте 42 на странице 9 и о токсичности во фрагменте 50 на странице 10. Оба не совпадают с автоматически выбранными эталонными страницами. И наоборот, фрагмент 154 на странице 36 касается отбора разметчиков, но не указывает место их набора. Пять первых фрагментов ограничены страницами 1--2 и не содержат ответов на три вопроса второй группы.

Последний прогон сохранил 12 ответов модели: на сокращенном индексе пять отказов и один неподтвержденный ответ о PPO-ptx; на полном - пять отказов и ответ «2.» со ссылкой на нерелевантный фрагмент. Ни один ответ не раскрывает нужный факт, поэтому средний балл по шкале 0--2 равен 0 для обоих индексов. Отказ при отсутствии контекста безопаснее догадки, но в этой шкале не является правильным ответом на вопрос.

В сокращенном индексе модель ответила: «PPO-ptx adds a variant of the Proximal Policy Optimization (PPO) algorithm to the original PPO training. [Фрагмент 1, стр. 2]». Ссылка неверна: фрагмент 1 относится к странице 1, а не 2. Ответ также не объясняет добавление обновлений предварительного обучения. Это пример ошибки генерации и оформления ссылки.

\subsection{Размер фрагментов}
Во всех вариантах перекрытие равно 100 символам. На данном наборе наибольшее совпадение страниц дали короткие фрагменты. Причину преимущества эти данные сами по себе не устанавливают. Это не универсальный оптимум: результат зависит от документа, вопросов и модели.
\begin{table}[H]\centering\caption{Размер фрагмента: ручная оценка top-3}\begin{tabular}{cccc}\toprule Размер&Фрагментов&Успешных из 10&Hit Rate@3\\\midrule 500&476&8&0.80\\1000&228&7&0.70\\2000&128&7&0.70\\\bottomrule\end{tabular}\end{table}
Здесь успех означает наличие хотя бы одного фрагмента, достаточного для ответа на вопрос из таблицы 1; это отличается от подсчета совпадающих страниц в выводе программы (6, 4 и 5). Для размера 500 вопрос о стадиях обучения не нашел нужный фрагмент, но ответ о PPO-ptx найден на стр. 9, хотя исходная автоматическая проверка ожидала стр. 4. Для вопроса о равных сравнениях ни один размер не вернул достаточный фрагмент. Оценка выполнена вручную по сохраненной выдаче и зависит от заданных эталонов.

\subsection{Перекрытие}
Размер зафиксирован как 1000 символов. Выбраны пять вопросов по границам на страницах 7, 9, 10, 18 и 37 при нулевом перекрытии. Критерий успеха - один фрагмент из top-3 содержит выбранный участок от 50 символов до границы до 50 символов после нее. Это проверяемый критерий сохранности участка, но его достаточноcть для ответа дополнительно оценивается по смыслу. Почти одинаковыми считались фрагменты с Jaccard-сходством множеств слов не менее 0.9. Загрузка PDF по страницам не позволяет перекрытию объединять текст через границу страниц.
\begin{table}[H]\centering\caption{Проверенные границы и перекрытие}\begin{tabular}{ccccc}\toprule Перекрытие&Фрагментов&Ответов из 5&Участков из 5&Повторов\\\midrule 0&217&3&0&0\\100&228&3&2&0\\200&246&5&4&0\\\bottomrule\end{tabular}\end{table}
Например, ответ о начальном составе очереди ожидания API разделен между двумя фрагментами при перекрытии 0. При перекрытии 100 или 200 найден единый фрагмент на странице 18, где указано, что первоначальные участники были сотрудниками OpenAI. Столбец «Ответов» получен ручной оценкой достаточности каждого найденного фрагмента. Столбец «Участков» показывает автоматический охват окна и отличается от содержательной оценки: например, ответ о reward model найден на странице 8 и без охвата окна на странице 7.

\subsection{Метод поиска}
Бинарный мешок слов использует косинусное сходство по словам из документа. По совпадению эталонных страниц он набрал больше баллов как для точных терминов, так и для перефразировок.
\begin{table}[H]\centering\caption{Сравнение поиска по содержанию фрагментов}\begin{tabular}{lccc}\toprule Метод&Точные, из 5&Перефразировки, из 5&По группам\\\midrule Мешок слов&4&1&0.80 / 0.20\\Плотные эмбеддинги&4&2&0.80 / 0.40\\\bottomrule\end{tabular}\end{table}
Для перефразированного вопроса о преобразовании GPT плотный поиск вернул фрагмент 2 со страницы 1 с описанием обучения на демонстрациях и обратной связи, хотя первоначальная эталонная страница - 3. Лексический поиск вернул страницы 12, 52 и 9. В исходном автоматическом выводе из-за совпадения страниц получены иные значения: 4/2 и 2/0. Семантическая оценка на пяти парах вопросов не доказывает универсальное превосходство метода.

\subsection{Проверяемые ссылки на источники}
Для проверки ссылок выполнены пять запросов на том же индексе; модель получала три найденных фрагмента, имя PDF, физическую страницу и номер фрагмента. Промпт требовал короткого ответа по контексту с точной копией метки источника. Модель 1B на первых пяти вопросах выдала четыре отказа и неполный ответ «2.» с нерелевантной ссылкой. Для отдельной проверки пяти ссылок использована модель 3B того же семейства и вопросы, для которых найден достаточный контекст. Наличие ссылки само по себе не подтверждает утверждение.

{\small\begin{longtable}{L{3.0cm}L{7.0cm}L{1.0cm}L{1.2cm}L{1.0cm}}
\caption{Проверка ответов и ссылок модели 3B}\\\toprule
Вопрос & Ответ модели & ID & Текст & Стр.\\\midrule\endfirsthead
\toprule Вопрос & Ответ модели & ID & Текст & Стр.\\\midrule\endhead
Стадии обучения & The three stages used to train InstructGPT are: (1) supervised fine-tuning, (2) reward model training, and (3) reinforcement learning via proximal policy optimization. [Фрагмент 9, стр. 3] & Да & Да & Да\\
Первые промпты & Because the project needed an initial source of instruction-like prompts to bootstrap the process. [Фрагмент 30, стр. 7] & Да & Да & Да\\
PPO-ptx & PPO-ptx adds pretraining updates to the PPO fine-tuning process to mitigate performance regressions on public NLP datasets. [Фрагмент 71, стр. 15] & Да & Да & Да\\
Категория запросов & The largest use-case category in the API prompt dataset was generation, with 45.6\% of the use-cases. [Фрагмент 25, стр. 6] & Да & Да & Да\\
Ограниченная группа ценностей & The learned alignment can reflect a limited group of values because it is based on the preferences of a specific human reference group for a specific application, rather than a broad consensus amongst many groups. [Фрагмент 89, стр. 18] & Да & Да & Да\\\bottomrule
\end{longtable}}
Столбец ID показывает наличие номера среди переданных фрагментов, «Текст» - подтверждение утверждения, «Стр.» - соответствие физической страницы метаданным. Промпт: Answer the question in one short sentence using only the context. End the answer by copying the exact [Фрагмент N, стр. M] tag at the beginning of the supporting passage. Figure numbers are not fragment numbers. Never create a tag. If the answer is absent, say Not enough information. При первоначальной формулировке модель перепутала номер рисунка 16 с номером фрагмента; поэтому в промпт добавлено указание не смешивать их.

\subsection{Отсутствие ответа в документе}
Пять вопросов без ответа: о точной годовой зарплате каждого разметчика, серийных номерах GPU, погоде в первый день обучения, точном счете за электроэнергию и MAC-адресе компьютера. Текст всех 68 страниц PDF проверен поиском соответствующих английских терминов. В статье есть упоминания оплаты и использования GPU, но требуемые точные значения не обнаружены. Вопрос о ссылке на репозиторий исключен: ссылка на репозиторий с образцами результатов есть на странице 10.

Для вопросов без ответа исходный промпт дал четыре отказа и один нерелевантный ответ: на вопрос о погоде модель вывела фрагмент о параметрах обучения. Строгий дал пять отказов. На вопрос о причинах первых промптов строгая модель отказалась, хотя фрагмент 30 на странице 7 содержал ответ. Исходная модель на вопрос о PPO-ptx зациклилась на повторении слова, не объяснив ответ.

Исходный промпт состоял из контекста и вопроса. Строгий предписывал отвечать только по переданным фрагментам, не выдумывать факты и ссылки, а при недостатке сведений выводить указанную фразу об отказе. В последнем прогоне обоим промптам передавались одни и те же заранее найденные фрагменты.

Полный шаблон исходного промпта: «Контекст: \{context\}; Вопрос: \{question\}; Ответ:». Строгий шаблон:
\begin{quote}\small
You are a retrieval QA system. Answer only with evidence present in the supplied PDF excerpts. If the excerpts do not contain enough evidence, output exactly: Not enough information in the document. Do not use outside knowledge. Do not invent a source, page, chunk identifier, or fact. For each factual sentence, append an existing citation exactly as shown in the excerpts: [Фрагмент N, стр. M]. Answer in English because the source document and questions are in English.

PDF excerpts: \{context\}. Question: \{question\}. Answer:
\end{quote}
\begin{table}[H]\centering\caption{Оценка ответов и отказов}\begin{tabular}{lccc}\toprule Промпт&Средний балл&Отказы из 5&Нерелевантных из 5\\\midrule Исходный&1.0&4&1\\Строгий&0.0&5&0\\\bottomrule\end{tabular}\end{table}
Средние баллы получены по пяти вопросам с ответом: для исходного промпта 2, 1, 0, 0, 2; для строгого 0, 0, 0, 0, 0. Оба промпта использовали одинаковые фрагменты для каждого вопроса.

\subsection{Уточняющие вопросы}
Реализован вариант Б: уточнение преобразуется в самостоятельный запрос до поиска. В трех диалогах хранится не более трех завершенных обменов, сессии независимы, очистка истории подтверждена выводом \texttt{[]}. Переформулирование может ошибаться: к вопросу об обучении InstructGPT модель добавила неподтвержденные подробности о данных.

Ниже приведены уточнения после первого вопроса каждой сессии. Запросы и страницы взяты из сохраненного вывода ноутбука; отказ обозначает ответ об отсутствии информации.
\begin{table}[H]\centering\small
\caption{Уточнения и поисковые запросы}
\begin{tabular}{L{3cm}L{5.3cm}L{1.4cm}L{1.6cm}L{1cm}}\toprule
Уточнение & Запрос для поиска & Стр. & Найден ответ & Балл\\\midrule
How was it trained? & How was InstructGPT trained on a massive dataset of text from the internet and books? & 7, 1, 2 & Часть & 0\\
Why was a reward model needed? & Why was a large language model with a reward structure needed to train InstructGPT? & 1, 9, 15 & Нет & 0\\
Which smaller model performed well? & Which smaller model was used for evaluation & 58, 56, 40 & Нет & 0\\
What did human raters prefer? & human raters preferred model 2 & 40, 58, 18 & Нет & 0\\
Where were they recruited from? & Where were the labelers recruited from? & 19, 36, 40 & Нет & 0\\
Why do their values matter? & Where were they recruited from? & 19, 36, 26 & Да & 0\\\bottomrule
\end{tabular}\end{table}
Первые вопросы диалогов: What is InstructGPT?, How were the models evaluated? и Who were the labelers?. На первый модель отказалась отвечать, на второй воспроизвела таблицу метрик с некорректными метками, на третий процитировала библиографическую запись вместо сведений о разметчиках. Первый поисковый запрос в таблице добавил неустановленные подробности; последний не сохранил смысл уточнения. В первом уточнении диалога обучения фрагмент 2, стр. 1 содержит частичное описание обучения; в последнем уточнении фрагмент 91, стр. 19 объясняет влияние разметчиков, однако модель ошибочно отнесла часть текста к другим фрагментам. В остальных четырех уточнениях достаточного фрагмента не найдено. Все шесть ответов получают 0 по шкале 0--2: пять являются отказами, последний содержит неверные ссылки.

В отдельном сравнении шесть уточнений выполнены и без преобразования запроса. Для How was it trained? исходный поиск вернул страницы 36, 56 и 19; после преобразования найдены страницы 7, 1 и 2, в том числе фрагмент 2 с описанием обучения. Это улучшило поиск, хотя модель все равно отказалась отвечать. Для Where were they recruited from? произошло обратное: исходный поиск нашел фрагмент 152 со страницы 36, где указаны Upwork и Scale AI, а после преобразования вернулся фрагмент 154 с той же страницы, не содержащий ответа. Преобразование запроса не всегда улучшает поиск.

\enlargethispage{3\baselineskip}
\section{Вывод}
Реализован RAG-конвейер и получены результаты семи экспериментов. При размере 500 и перекрытии 100 символов ручная оценка дала 8 успешных поисков из 10. При сравнении двух методов на пяти парах вопросов результаты зависели от формулировки; объем выборки не позволяет обобщать вывод. На пяти выбранных границах сохранность участков в top-3 увеличилась при большем перекрытии. Строгий промпт нередко отказывался даже при наличии ответа в контексте.
\end{document}
